\section{Schema/Interface Sensitivity}
\label{app:schema-interface-sensitivity}

The current artifact uses the clarified schema for all main results. Because field definitions can affect whether evidence roles are expressed as support, rejected evidence, or verification actions, we ran a small schema-role decomposition ablation over generated-lore and buried-primary rows after the main role-uncued pilot. The ablation holds the evidence packet fixed and varies only the structured output interface. It is exploratory follow-up evidence, not a replacement for the main pilot table.

The interface finding that remains valid in this pilot is narrow. Structured fields make evidence and action allocation inspectable, but they are not the paper's main research object. In generated-lore rows, strict failures are sensitive to whether the scorer treats polluted documents in support fields as accepted support or recognizes that prose and actions diagnosed them as unverified. In active-verification rows, the action schema makes executable next steps measurable instead of leaving them as prose.

A broader follow-up schema rerun used 36 selected visible-view source tasks: 12 generated-lore tasks and 8 tasks each from buried-primary, conflicting-evidence, and false-consensus conditions. With two models and two schema variants, this produced 144 schema-rerun rows and 72 paired current-versus-role-decomposed comparisons. The role-decomposed schema changed strict contract success only slightly on average (mean delta 0.014) and did not create a new model ranking. It did, however, make the contract easier to inspect by separating clean support, refuting evidence, rejected or contaminated evidence, diagnostic evidence, and verification actions. The rerun was not uniformly better on every component: in this selected slice, clean-support recovery was lower under the role-decomposed schema (0.222 versus 0.667), even while mean strict contract success rose slightly. This should be read as an interface-sensitivity appendix, not as a central claim about model cognition or as an explanation for every main-pilot failure.
